\chapter{Gli SDK multi-linguaggio}
\label{cap:18-sdk}

\nf{} fornisce SDK per sei ambienti: Java, Java-lite, Python, Go, JavaScript e
un adattatore per eseguibili generici (\code{exec}/bash). Questo capitolo li
attraversa e --- pi\`u interessante --- descrive il meccanismo con cui il
progetto verifica che parlino tutti lo stesso contratto.

\section{Il contratto minimo}

Un SDK di funzione ha tre obblighi: consentire all'autore di dichiarare un
handler, esporlo secondo uno dei tre protocolli di runtime
(\cref{cap:17-runtime-funzione}), e tradurre il valore restituito in una
risposta conforme.

L'interfaccia lato Java \`e un solo metodo:

\begin{lstlisting}[language=Java]
public interface FunctionHandler {
    Object handle(InvocationRequest request);
}
\end{lstlisting}

Ogni SDK realizza questo contratto con l'idioma del proprio linguaggio.

\section{Java: due SDK, due filosofie}

\subsection{\code{sdks/java}}

L'SDK completo (1\,607 righe) \`e costruito su Spring Boot. La dichiarazione di
un handler \`e un'annotazione:

\begin{lstlisting}[language=Java]
/**
 * Marks a Spring bean as a Nanofaas function handler.
 *
 * <p>Spring component scanning, not user code, discovers this annotation and registers the
 * handler in the application context. The runtime then resolves the single active handler bean
 * from that context when the control plane calls {@code /invoke}.</p>
 */
@Target(ElementType.TYPE) @Retention(RetentionPolicy.RUNTIME) @Documented @Component
public @interface NanofaasFunction {
    String value() default "";
}
\end{lstlisting}

L'autore della funzione ottiene iniezione delle dipendenze, configurazione, e
l'intero ecosistema Spring. Paga in impronta di memoria e tempo di avvio.

\subsection{\code{sdks/java-lite}}

L'SDK leggero (1\,454 righe) rinuncia a Spring. Costruisce direttamente sul
\code{HttpServer} del JDK, con tre gestori --- \code{InvokeHandler},
\code{HealthHandler}, \code{MetricsHandler} --- e nessun contenitore di
inversione del controllo.

Il confronto fra i due, misurato sulla stessa funzione
(\cref{cap:24-footprint}), \`e la giustificazione della loro coesistenza:

\begin{center}
\small
\begin{adjustbox}{max width=\textwidth}
\begin{tabular}{@{}lrr@{}}
\toprule
 & \textbf{\code{word-stats-java}} & \textbf{\code{word-stats-java-lite}} \\
\midrule
build & Spring Boot su JVM & binario nativo GraalVM \\
SDK & \code{sdks/java} & \code{sdks/java-lite} \\
residente, a riposo & 133\,MiB & 2{,}7\,MiB \\
residente, a caldo & 160\,MiB & 51{,}6\,MiB \\
tempo di servizio & 5{,}57\,ms & 2{,}15\,ms \\
\bottomrule
\end{tabular}
\end{adjustbox}
\end{center}

Il fattore cinquanta a riposo \`e la differenza fra un framework e la sua
assenza; il fattore tre a caldo mostra che il carico di lavoro effettivo
riavvicina i due (\cref{cap:24-footprint}).

\begin{quote}
\textbf{Avvertenza sperimentale.} La documentazione di progetto registra che i
commenti degli scenari hanno per un certo tempo trattato le due funzioni
\code{word-stats} come intercambiabili. Non lo sono: differiscono per SDK, per
tipo di build e per tempo di servizio. \emph{Qualunque esperimento che le tratti
come equivalenti \`e confuso.} \`E il tipo di errore che una piattaforma con
molte varianti dello stesso artefatto rende facile commettere, ed \`e riportato
qui perch\'e la sua scoperta ha invalidato analisi precedenti.
\end{quote}

\section{Python}

L'SDK Python \`e costruito su FastAPI (1\,407 righe di sorgente). La
dichiarazione dell'handler \`e un decoratore:

\begin{lstlisting}[language=Python]
@nanofaas_function
def handle(input):
    return {"echo": input}
\end{lstlisting}

Il decoratore registra la funzione in una variabile a livello di modulo, con la
regola che un solo handler pu\`o essere attivo e che una seconda applicazione
sovrascrive la prima. Sono supportati sia handler sincroni sia \code{async}.

La scelta di una registrazione globale anzich\'e di un registro esplicito \`e la
forma pi\`u idiomatica in Python e la meno cerimoniosa; ha il difetto di rendere
l'ordine degli import significativo, che \`e il prezzo abituale di questa
tecnica.

\section{Go, JavaScript, exec}

L'SDK Go (1\,914 righe) incorpora un runtime HTTP nel binario della
funzione ed espone un proprio endpoint \code{/metrics} Prometheus, con contatori
lato runtime per invocazioni, durata dell'handler, cold start e callback
asincrone scartate. \`E l'unico SDK che strumenta autonomamente il proprio lato
del confine, il che lo rende lo strumento pi\`u diretto per separare il tempo
speso dentro l'handler da quello speso nel trasporto.

L'SDK JavaScript (1\,531 righe di TypeScript) espone un runtime HTTP analogo per Node.

L'adattatore \code{exec} non \`e un SDK ma il modo \code{STDIO} del watchdog: un
eseguibile che legge JSON da stdin e scrive JSON su stdout \`e una funzione
valida senza libreria alcuna. \`E la via d'accesso per gli script bash e per i
binari preesistenti.

\section{Il contratto di saturazione}
\label{sec:18-saturazione}

Un runtime che accetta lavoro senza un limite \`e un runtime che, sotto carico,
consuma la memoria del contenitore finch\'e il kernel non lo termina. Il progetto
ha quindi fissato, in un unico documento (\code{sdks/runtime-contract/README.md}),
la politica con cui \emph{ogni} runtime di funzione deve rifiutare, e ne ha
fatto un contratto verificabile allo stesso modo della parit\`a funzionale.

Il contratto vale sia per le chiamate dirette a \code{POST /invoke} sia per i
dispacci del piano di controllo. L'ammissione \`e \emph{fail-fast}: il runtime
riserva capacit\`a di conteggio e di byte per ingresso, handler e callback
\emph{prima} di avviare l'handler o di produrre una copia serializzata di grandi
dimensioni. Una callback richiesta dall'invocazione fa parte dell'ammissione: un
runtime non pu\`o accettare il lavoro e scartare in silenzio la sua callback
terminale.

\begin{center}
\small
\begin{adjustbox}{max width=\textwidth}
\begin{tabular}{@{}llp{6.6cm}@{}}
\toprule
\textbf{Esito} & \textbf{Codice} & \textbf{Quando} \\
\midrule
429 & \code{RUNTIME\_HANDLER\_SATURATED} & ammissione dell'handler satura \\
429 & \code{RUNTIME\_CALLBACK\_SATURATED} & capacit\`a di callback satura \\
413 & \code{RUNTIME\_INPUT\_TOO\_LARGE} & ingresso oltre il limite di byte \\
500 & \code{RUNTIME\_OUTPUT\_TOO\_LARGE} & uscita scoperta \emph{dopo} l'esecuzione;
l'handler \`e partito, l'uscita non viene trattenuta e una callback d'errore
limitata \`e comunque dovuta \\
503 & \code{RUNTIME\_STOPPING} & richiesta accettata dopo l'inizio dell'arresto \\
504 & \code{HANDLER\_TIMEOUT} & attesa dell'handler scaduta, a livello di
tentativo (non il \code{408} per singolo attendente) \\
\bottomrule
\end{tabular}
\end{adjustbox}
\end{center}

Gli errori hanno forma \code{\{"error":\{"code":...,"message":...\}\}} e quelli
ritentabili portano \code{Retry-After: 1}. Tre regole completano la politica.
Ogni attesa \`e finita e configurabile. Ogni riserva di conteggio e di byte resta
posseduta finch\'e il lavoro rappresentato e il carico trattenuto non sono
\emph{fisicamente} spariti: scadere l'attesa HTTP non prova che un handler non
cooperativo si sia fermato, e per questo il runtime Go, per esempio, documenta
che un handler scaduto continua a occupare il proprio posto. Infine i runtime non
ridispacciano mai: la decisione di ritentare appartiene al piano di controllo
(\cref{cap:06-nucleo}), che la esercita mantenendo stabile
\code{X-Execution-Id} e incrementando \code{X-Dispatch-Attempt}.

Il contratto \`e eseguibile. Le casistiche sono in un solo file,
\code{saturation-wire-corpus.json}, e ciascuna implementazione ha un adattatore
di test sottile che lo legge; un validatore Python
(\code{validate\_saturation\_wire\_corpus.py}) verifica la coerenza interna del
corpus --- riferimenti, azioni e barriere, esiti esatti, contatori finali a
regime --- e ha una propria suite di mutazioni, cos\`i che il corpus non possa
diventare un oracolo indipendente e incoerente. L'ordinamento fra azioni usa
barriere nominate, non attese temporizzate: il tempo trascorso non \`e evidenza
di ordine. Gli adattatori sono nella suite nativa di ciascun linguaggio; Java e
Java-lite deserializzano il modello tipizzato completo, Go rifiuta i campi
sconosciuti.

I limiti hanno valori predefiniti finiti e sono configurabili per variabile
d'ambiente. Il runtime Go e quello Python condividono lo stesso vocabolario
(\code{NANOFAAS\_MAX\_CONCURRENT\_HANDLERS}, 32 predefiniti;
\code{NANOFAAS\_MAX\_INPUT\_BYTES} e \code{NANOFAAS\_MAX\_OUTPUT\_BYTES}, 1\,MiB;
\code{NANOFAAS\_MAX\_PENDING\_CALLBACKS}, 128;
\code{NANOFAAS\_MAX\_CALLBACK\_PAYLOAD\_BYTES}, 2\,MiB;
\code{NANOFAAS\_MAX\_PENDING\_CALLBACK\_BYTES}, 16\,MiB; timeout di lettura del
corpo, di tentativo di callback e di arresto di 5\,s), e valori oltre i massimi
di produzione ricadono sui predefiniti prima che venga allocato alcunch\'e. Un
dettaglio che mostra il criterio: la capacit\`a di callback effettiva \`e il
minimo fra il limite di conteggio e il rapporto fra il tetto aggregato di byte e
il tetto per singola callback, perch\'e il runtime riserva un'intera quota di
callback prima di avviare l'handler.

\`E la stessa filosofia dell'ArchUnit del piano di controllo
(\cref{sec:26-archunit}) applicata ai runtime: un limite che non \`e imposto e
verificato \`e un'intenzione. Il divario iniziale era reale ---
l'inventario del documento registra, per ciascun SDK, code e attivit\`a
\emph{senza} un tetto di ammissione (Java e Java-lite con un thread virtuale per
richiesta, Python con la coda illimitata di \code{to\_thread}, Go con una
goroutine per richiesta) --- e la sua chiusura \`e avvenuta per fasi che il
documento numera (P16a fissa politica e inventario, P16b implementa i limiti su
tutti i runtime, P17 il tracciamento del lavoro fisico in Python, P18 le risorse
di arresto di Java-lite). Non \`e stata misurata qui l'impronta di questi limiti
sul percorso caldo.

\section{La parit\`a fra SDK, e come viene verificata}

\subsection{Il problema}

Sei implementazioni indipendenti dello stesso contratto divergono. Non
divergono in modo spettacolare --- ciascuna \`e testata --- ma nei dettagli:
come viene serializzato un valore nullo, che cosa accade con un input vuoto, come
viene rappresentato un errore. Divergenze di questo tipo sono invisibili finch\'e
qualcuno non confronta due SDK sulla stessa funzione, che \`e esattamente ci\`o
che un esperimento comparativo fa.

\subsection{La soluzione: fixture condivise}

Il progetto definisce tre \emph{famiglie} di funzioni di riferimento ---
\code{word-stats}, \code{json-transform}, \code{roman-numeral} --- e per ciascuna
un file di fixture JSON con input e output attesi. Ogni implementazione consuma
lo stesso file dalla propria suite di test nativa.

\begin{center}
\small
\begin{adjustbox}{max width=\textwidth}
\begin{tabular}{@{}lccc@{}}
\toprule
\textbf{SDK} & \textbf{\code{word-stats}} & \textbf{\code{json-transform}} &
\textbf{\code{roman-numeral}} \\
\midrule
Java & s\`i & s\`i & s\`i \\
Java Lite & s\`i & s\`i & s\`i \\
Go & s\`i & s\`i & s\`i \\
Python & s\`i & s\`i & s\`i \\
JavaScript & s\`i & s\`i & s\`i \\
exec/bash & s\`i & s\`i & s\`i \\
\bottomrule
\end{tabular}
\end{adjustbox}
\end{center}

Il cancello completo si esegue con un comando:

\begin{lstlisting}[language=bash]
./functions/contract-tests/run.sh
\end{lstlisting}

Diciotto funzioni, sei SDK, tre famiglie. \`E il meccanismo che converte
un'affermazione di parit\`a in una verifica.

\subsection{Fixture di correttezza e fixture di carico}

Il progetto distingue due tipi di fixture, e la distinzione \`e metodologicamente
significativa:

\begin{itemize}[leftmargin=*]
  \item \code{correctness.json} contiene input \emph{e} output attesi. Verifica
  la semantica.
  \item \code{performance-*.json} contiene input soltanto, e \emph{omette
  deliberatamente} gli output attesi.
\end{itemize}

L'omissione \`e voluta. Una fixture di carico serve a fissare l'input dei
benchmark cos\`i che due esecuzioni siano confrontabili; verificare l'output
durante un benchmark aggiungerebbe alla misura il costo del confronto e
introdurrebbe una variabile nel percorso caldo. Separare i due ruoli evita che
uno strumento serva a due scopi facendone bene nessuno.

\section{Il catalogo delle funzioni di esempio}

Il repository contiene funzioni di esempio in cinque linguaggi:

\begin{center}
\small
\begin{adjustbox}{max width=\textwidth}
\begin{tabular}{@{}lp{8.4cm}@{}}
\toprule
\textbf{Linguaggio} & \textbf{Funzioni} \\
\midrule
Java & \code{word-stats}, \code{json-transform}, \code{roman-numeral},
\code{figlet}, \code{qr-code}, \code{handler-envelope},
\code{binary-envelope}, pi\`u tre varianti \code{-lite} \\
Python & le stesse tre famiglie, pi\`u \code{qr-code},
\code{handler-envelope}, \code{mlimage} \\
Go, JavaScript & le tre famiglie, pi\`u \code{qr-code} e
\code{handler-envelope} \\
bash & le stesse cinque, pi\`u \code{figlet} \\
\bottomrule
\end{tabular}
\end{adjustbox}
\end{center}

Le funzioni non sono soltanto dimostrative. Ciascuna copre un aspetto del
contratto: \code{word-stats} \`e il carico di riferimento delle campagne di
misura; \code{handler-envelope} esercita la busta di risposta
(\cref{cap:17-runtime-funzione}); \code{binary-envelope} esercita la codifica
base64; \code{qr-code} e \code{figlet} producono output non testuali;
\code{mlimage} \`e l'unica con un tempo di esecuzione sostanziale e un'immagine
pesante, e serve a esercitare il cold start reale.
